\section{Cálculo das Probablidades}

Para espaços amostrais enumeráveis, ao associar a cada resultado individual
um número em $[0, 1]$, então a soma dos pesos dos resultados presentes num evento
é a probablidade desse evento.

\begin{thm:soma prob}
\label{thm:soma prob}
Seja $\Omega = \{\omega_i\}_{\mathbb{N}}$ um espaço amostral enumerável, e $\mathcal{B}$ uma
$\sigma$-Álgebra de $\Omega$.
Seja $(p_i)_\mathbb{N}$ uma sequência de números não negativos cuja a soma é $1$.
Presuma também que, se $A \in \mathcal{B}$, então
\[
	P(A) = \sum_{w_i \in A} p_i.
\]
A tripla $(\Omega, \mathcal{B}, P)$ é um espaço de probabilidade.
\begin{proof}
	Precisamos provar que $P$ é uma medida de probabilidade que satisfaz os Axiomas de
	Kolmogorov.

	\paragraph{Não negatividade} a soma de números não negativos é não negativa, então
	para todo $A$ em $\mathcal{B}$ vale que $P(A) \geq 0$.

	\paragraph{Probablidade do espaço amostral} A série $\sum_{i=1}^\infty$ deve coicidir
	com a probabilidade de $\Omega$, série essa que será igual a $1$ por hipotese.

	\paragraph{Soma de eventos disjuntos} Seja $\{A_j\}_{\mathbb{N}}$ uma família
	disjunta de eventos de $\mathcal{B}$ tal que
	\[
		A_j = \{\omega_i \mid i \in I_j\}
	\]
	em que $I_j$ é um conjunto de índices. Tão logo obtemos
	\[
		P\left(\bigcup_{j=1}^\infty A_j\right) = \sum_{j=1}^\infty \sum_{i \in I_j} p_{ji}
		= \sum_{j=1}^\infty P(A_j)
	\]
\end{proof}
\end{thm:soma prob}

Esse resultado é útil para computar a probablidade de uma vasta gama de
experimentos aleatórios na prática. Um exemplo é o lançamento de um dado de 6
faces; se acreditamos que o dado é honesto, isto é, que o dado não é viciado em
nenhuma de suas faces, então podemos associar $1/6$ a cada face do dado. Se por
exemplo, deseja-se saber a probabilidade de obtermos um número menor ou igual a
3, basta somar as faces que satisfazem essa condição, obtendo uma probablidade
de $0.5$.
Experimentos os quais acreditamos não conterem vícios são chamados de equiprováveis.

As $\sigma$-Álgebras que equivalem ao conjunto das partes do espaço amostral são
chamadas de $\sigma$-Álgebra potência.
Quando o espaço amostral do experimento é enumerável, ocorre que cada subconjunto unitário
do espaço amostral -- um resultado individual -- é um evento do experimento. Se o
experimento é equiprovável, então todos os resultados são eventos de mesma probablidade.

\begin{def:evento equiprovavel}
Seja o espaço $(\Omega, \mathcal{B}, P)$ com $\Omega$ finito, $\mathcal{B} = \Parts{\Omega}$
e $P$ definido como no Teorema \ref{thm:soma prob}.
Chama-se esse espaço de equiprovável se, e somente se,
\[
	\forall i \in [n], P(\omega_i) = \frac{1}{n}
\]
onde $n = |\Omega|$.
\end{def:evento equiprovavel}

A probablidade de eventos em experimentos com resultados equiprováveis destaca-se
por ser intuitiva de calcular. A probablidade de evento será, simplesmente, a
fração dos resultados presentes nele. Esse resultado será enunciado como Corolário
do Teorema \ref{thm:soma prob}.

\begin{cor:equiprovavel}
\label{prop:equiprovavel}
A probabilidade de um evento $A \in \Parts{\Omega}$ num espaço $(\Omega, \Parts{\Omega}, P)$
equiprovável é igual a
\[
	\frac{|A|}{|\Omega|}
\]
\begin{proof}
	\[
		P(A) = |A| \cdot \frac{1}{|\Omega|} = \frac{|A|}{|\Omega|}
	\]
\end{proof}
\end{cor:equiprovavel}
